% ============================================================
% CONCLUSIONES — TFG
% Índice Sintético de Turismo Sostenible por CCAA (2016–2024)
% ============================================================

\chapter{Conclusiones}

\section{Hallazgos principales}

Este trabajo ha construido un índice sintético de sostenibilidad turística para las 19 comunidades autónomas españolas (2016--2024), comparando cinco combinaciones metodológicas. Los hallazgos principales son:

\subsection{1. Liderazgo robusto del norte}

Navarra encabeza el ranking en los años 2022, 2024 (Casos 1, 2, 4) y posición 3 en Caso 5. Su desviación estándar entre metodologías es $\sigma \approx 1.0$, muy baja comparada con otras CCAA. Asturias, Galicia y La Rioja ocupan consistentemente posiciones 2--5. Este resultado es **robusto**: el liderazgo norte-noroeste persiste bajo prácticamente todas las variantes metodológicas, lo que indica que no es artefacto sino reflejo de estructura real de datos.

\subsection{2. Equivalencia Caso 1 vs Caso 3}

La comparación directa entre normalización min-max (Caso 1) y z-score (Caso 3) bajo agregación lineal produce **rankings idénticos** (correlación Spearman $\rho_s = 1.0$). Esto confirma que, bajo ponderación PCA lineal, la elección de escala de normalización no afecta al orden relativo. El debate min-max vs. z-score es académico en este contexto.

\subsection{3. Función de agregación determina más que normalización}

Comparando Caso 1 (media aritmética) con Caso 2 (media geométrica), ambos bajo min-max y PCA(PC1), se observan cambios sustanciales. Baleares desciende de posición 11 a 10 con score dramáticamente reducido (46.94 → 22.64). Esto refleja que **media geométrica penaliza desequilibrios** entre dimensiones: Baleares es fuerte en ECO pero débil en ENV/SOC. Bajo media aritmética, compensación; bajo media geométrica, penalización.

\subsection{4. Caso 4 captura varianza residual}

El uso de todas las componentes principales (Caso 4) en lugar de solo PC1 amplifica la importancia de capas de variación secundarias. Asturias y Galicia emergen más fuertes porque su fortaleza ambiental se amplifica. La dimensión económica (α=0.4711) es ligeramente superior a Social (β=0.3950) y Ambiental (δ=0.4065), pero el peso de cada una está cerca del tercio (normalizados: 37%, 31%, 32%).

\subsection{5. Caso 5 revela estructura multifactorial}

La rotación Varimax desglosa ECO en 7 factores independientes, SOC en 4, ENV en 2. Esto demuestra que ``sostenibilidad económica'' del turismo regional no es unitaria sino agregado multidimensional: tarificación, rentabilidad, ocupación, gasto, inflación, empleo y variaciones operan con lógica propia. Galicia lidera bajo Caso 5 porque su balance entre estos 7 factores es óptimo, no porque super-destaque en ninguno.

\subsection{6. 2020 como perturbación estadística}

La pandemia anuló varianza en 11 indicadores en 2020, homogeneizando artificialmente el panel. Incluir 2020 afecta los pesos PCA de forma no trivial. El análisis posterior (2021--2024) muestra valores más altos e diferenciación reaparece, confirmando que 2020 fue evento transversal, no revelador de debilidades estructurales.

\section{Limitaciones del análisis}

\subsection{Cobertura ambiental sesgada}

Dimensión ENV tiene solo 4 indicadores activos tras depuración (versus 20--25 en ECO, 15--20 en SOC). Dos son complementarios (IAMB0006, IAMB0007: composición modal transporte). Aspectos ambientales relevantes (huella carbono del alojamiento, gestión de residuos, presión sobre espacios naturales) no están representados. El peso de ENV (45.17\%) oculta poca diversidad: dimensión modal-céntrica, no ambiental en sentido amplio.

\subsection{Heterogeneidad de disponibilidad}

Eliminación de indicadores con disponibilidad $< 20\%$ fue necesaria pero arbitraria. Muchos restantes tienen 20--30\%. Imputación MICE asume MAR; si patrón es MNAR (ej. CCAA no reporta variable porque no es relevante), sesgo de relleno.

\subsection{Multicolinealidad en ECO no resuelta por PCA simple}

Casos 1--3 (PCA simple PC1) no resuelven redundancia de 7 indicadores económicos con $r > 0.95$. PC1 amplifica corr pesos a los demás. Casos 4 y 5 atacan esto mejor: Caso 4 usa todas las componentes; Caso 5 usa rotación Varimax.

\subsection{Ausencia de denominadores poblacionales}

Indicadores de volumen no están relativizados por población o superficie. Sesgo que favorece CCAA pequeñas con mucho turismo relativo (Baleares aún lidera en absoluto pero su per-cápita sería más alto).

\section{Problemas técnicos encontrados}

\subsection{Media geométrica ponderada: implementación manual}

Funciones estándar COINr para media geométrica generaban errores. Solución manual:

\begin{equation}
    G = \exp\left(\sum_{j=1}^{J} w_j^* \ln x'_j\right)
\end{equation}

requirió garantizar $x > 0$ (rango [1, 100]). Vulnerabilidad computacional.

\subsection{Winsorización año a año pierde perspectiva histórica}

Tratamiento coin-por-coin produce perda de información temporal. Valor extremo en 2016 se trata independiente de si es normal en 2024.

\section{Elección del método recomendado}

\textbf{Se recomienda el Caso 5 (PCA/FA rotado OCDE + Media Aritmética) como índice oficial.}

\subsection{Justificación detallada}

\begin{enumerate}
    \item \textbf{Resolución de multicolinealidad}: Rotación Varimax produce loadings concentrados. Indicadores redundantes quedan en mismo factor, evitando sobrerepresentación. Caso 5 identifica ECO como 7 factores independientes, no uno. La sostenibilidad económica del turismo regional es inherentemente multidimensional.
    
    \item \textbf{Adhesión al estándar OCDE}: El manual de \emph{Composite Indicators} \citep{oecd2008handbook} establece explícitamente rotación Varimax como Paso 3 post-depuración. Caso 5 implementa metodología de referencia en construcción de indicadores públicos europeos. Comparabilidad internacional.
    
    \item \textbf{Mejor cuadra de balances dimensionales}: Bajo PCA simple (Casos 1--3), si PC1 de ECO explica 48\% versus 40--41\% en ENV/SOC, PC1-Eco domina automáticamente. Rotación Varimax del Caso 5 equilibra distribuyendo pesos entre factores ortogonales.
    
    \item \textbf{Interpretabilidad mejorada}: Los 7 factores económicos tienen nombres conceptuales (tarificación, rentabilidad, ocupación, gasto, inflación, empleo). Facilita comunicación a gestores públicos. Un gestor comprende ``sostenibilidad económica depende de 7 dinámicas'' más que ``PC1 explica 48\% varianza''.
    
    \item \textbf{Robustez ante cambios de especificación}: Rotación hace ranking menos sensible a inclusión/exclusión de indicador correlacionado. En PCA simple, eliminar IECO0033 cambiaría loadings de PC1 dramáticamente. En Caso 5, lo desplazaría dentro del factor rotado con impacto menor.
    
    \item \textbf{Penalización de super-especialización}: Galicia lidera en Caso 5 porque su equilibrio multifactorial es robusto, no por un único indicador dominante. Baleares, pese a fuerza económica, desciende a posición 15 porque sus demás dimensiones son débiles. Esto alinea el índice con concepto de sostenibilidad multidimensional.
\end{enumerate}

\subsection{Comparación vs resto de casos}

\textbf{Debilidades de Casos 1--4}:

\begin{itemize}
    \item \textbf{Caso 1}: Simple PC1 no resuelve multicolinealidad. ECO domina artificialmente.
    \item \textbf{Caso 2}: Media geométrica es conceptualmente atractiva para penalizar desequilibrios, pero no resuelve multicolinealidad. Además no está en manual OCDE.
    \item \textbf{Caso 3}: Equivalente a Caso 1, ranking idéntico. Innecesario.
    \item \textbf{Caso 4}: Usa todas las componentes, mejora sobre Casos 1--3, pero fórmula STI con α, β, δ como exponentes es más opaca. Media geométrica ponderada requiere cálculo manual propenso a errores.
\end{itemize}

\subsection{No existe decisión objetivamente óptima}

La recomendación de Caso 5 se basa en criterios de robustez estadística, adherencia a estándares, interpretabilidad. **Pero** depende del contexto:

\begin{itemize}
    \item Si deseas máxima **penalización de desequilibrios**: Caso 2 (media geométrica)
    \item Si deseas capturar **residual variability**: Caso 4
    \item Si deseas máxima **simplicidad comunicativa**: Caso 1
\end{itemize}

La contribución más honesta: el trabajo documentó que **las conclusiones dependen de decisiones metodológicas**. Que esas decisiones sean **explícitas y transparentes** ante usuarios finales (tomadores de decisión, ciudadanía).

\section{Líneas de trabajo futuro}

\subsection{Enriquecimiento de dimensión ambiental}

Incorporar:
- Huella carbono por pernoctación (Ministerio Transición Ecológica)
- Consumo agua por pernoctación (servicios municipales)
- Certificación ambiental de establecimientos
- Presión sobre espacios protegidos

Una ENV ampliada (12--15 indicadores) reduciría sesgo por transporte.

\subsection{Análisis de sensibilidad formal Monte Carlo}

Según manual OCDE, variar simultáneamente umbrales de Skewness/Kurtosis, número de PCs, cutoffs de correlación. Calcular intervalos de confianza para posiciones. Identificar CCAA con posiciones ``sólidas'' vs. ``frágiles''.

\subsection{Tratamiento diferenciado de 2020}

- Test de Chow para detectar ruptura estructural
- Variable ficticia en PCA
- Análisis por subperíodos (pre-pandemia, pandemia, post-pandemia)

\subsection{Desagregación a nivel provincial (NUTS~3)}

Actual NUTS~2 oculta heterogeneidad intra-regional. Málaga vs. Córdoba dentro Andalucía son muy distintos. Extensión provincial requiere resolver disponibilidad de datos.

\subsection{Indicadores emergentes}

- Overtourism (concentración geográfica vs población)
- Resiliencia a crisis (diversificación de mercados)
- Equidad local (retención de ingresos vs cadenas globales)

\subsection{Comparación con marcos internacionales}

- ETIS (European Tourism Indicators System)
- UNWTO Sustainable Tourism Index
- Foro Económico Mundial Travel & Tourism Competitiveness

Validación externa contra índices establecidos.

\section{Síntesis conclusiva}

Este trabajo proporciona **instrumento reproducible y riguroso** para monitorear sostenibilidad turística regional española.

El índice no resuelve automáticamente la pregunta política «¿Qué constituye verdadero turismo sostenible?» —eso requiere deliberación multi-stakeholder—. Pero proporciona **evidencia cuantitativa rigurosa** sobre posiciones relativas comparables temporalmente.

Es herramienta de **diagnóstico y monitoreo**, no prescripción. Debe usarse junto con análisis cualitativo territorial, encuestas de comunidades locales, marco regulatorio que traduzca indicadores en política pública.

La comparación de cinco casos metodológicos demuestra que robustez es **parcial**:
- **Consenso**: norte lidera bajo casi todas las alternativas = sostenibilidad real
- **Ambigüedad**: Baleares, Madrid, regiones especializadas dependen fuertemente de metodología = necesario explicitar qué tipo de sostenibilidad se prioriza

Futuras ediciones que incorporen los elementos sugeridos (ENV ampliada, análisis de sensibilidad formal, tratamiento de 2020) incrementarán confianza aplicable a toma de decisión pública.

\end{document}
